\section{Comparação entre timeout fixo e adaptativo}

\preencher{Escolher GBN ou SR e registrar a justificativa da escolha. Executar
o mesmo protocolo, com a mesma janela, arquivo, semente e condições, usando
timeout fixo e adaptativo sob perda de 5\% e 10\%.}

\begin{table}[H]
  \centering
  \caption{Comparação entre timeout fixo e adaptativo.}
  \label{tab:timeout}
  \begin{tabular}{llrrrr}
    \toprule
    Protocolo & Perda & Timeout & Retransmissões & Tempo total (s) &
      Throughput médio (kbps) \\
    \midrule
    \preencher{GBN/SR} & 5\%  & fixo       & \preencher{n} & \preencher{s} & \preencher{kbps}\\
    \preencher{GBN/SR} & 5\%  & adaptativo & \preencher{n} & \preencher{s} & \preencher{kbps}\\
    \preencher{GBN/SR} & 10\% & fixo       & \preencher{n} & \preencher{s} & \preencher{kbps}\\
    \preencher{GBN/SR} & 10\% & adaptativo & \preencher{n} & \preencher{s} & \preencher{kbps}\\
    \bottomrule
  \end{tabular}
\end{table}

\subsection{Interpretação}
\preencher{Explicar em que situação o timeout adaptativo tem maior vantagem e
em que situação o timeout fixo é suficiente. Relacionar os resultados às
perdas, ao RTT e às retransmissões espúrias.}
